% FILE: 04_implementation_surface.tex
% Project: otel-weaver-exp-002-diff-to-residual
% Thesis Role: weaver-diff-to-pi-residual-bridge-experiment

\section{Implementation Surface}
\label{sec:otel-weaver-exp-002-diff-to-residual-impl}

\subsection{Source Files}
\label{sec:otel-weaver-exp-002-diff-to-residual-impl-sources}

The experiment's implementation surface consists of two Rust files and one
manifest, all located under the experiment root at
\texttt{otel-weaver/experiments/exp-002-weaver-diff-to-pi-residual/}:

\paragraph{src/bridge.rs}
The sole implementation file. It declares four public items:
\begin{enumerate}
    \item \texttt{TelemetryFeedstockS2} --- derives \texttt{Debug}, \texttt{Clone},
          \texttt{Serialize}, \texttt{Deserialize}; carries \texttt{instance\_id},
          \texttt{transition\_name}, \texttt{token\_color: Option<String>}, and
          \texttt{witness\_hash}. Serde rename attributes map Rust field names to
          their OTel dotted-path attribute names.
    \item \texttt{TelemetryFeedstockS1} --- derives the same traits; carries
          \texttt{instance\_id}, \texttt{activity\_name}, and \texttt{witness\_hash}.
    \item \texttt{BridgeRx} --- does not derive \texttt{Serialize}/\texttt{Deserialize}
          (it is a processing struct, not a data struct); holds
          \texttt{pub rename\_map: HashMap<String, String>}; implements \texttt{new()}
          and \texttt{translate()}.
    \item The \texttt{\#[cfg(test)] mod tests} block containing
          \texttt{test\_bridge\_translation\_equivalence}.
\end{enumerate}

The file imports \texttt{std::collections::HashMap} and
\texttt{serde::\{Serialize, Deserialize\}}. No other internal crate modules are
referenced. The library target is the entire file, as declared in \texttt{Cargo.toml}
with \texttt{path = "src/bridge.rs"}.

\paragraph{Cargo.toml}
Package manifest declaring:
\begin{verbatim}
[package]
name = "exp-002-weaver-diff-to-pi-residual"
version = "0.1.0"
edition = "2021"

[lib]
name = "bridge"
path = "src/bridge.rs"

[dependencies]
serde = { version = "1.0", features = ["derive"] }
serde_json = "1.0"
serde_yaml = "0.9"
\end{verbatim}

Note that \texttt{serde\_json} and \texttt{serde\_yaml} are declared as
dependencies but are not imported in \texttt{src/bridge.rs}. They are present to
support future serialization of feedstock records from JSON and YAML sources at the
ingestion boundary. The \texttt{Cargo.lock} confirms that dependency resolution
was completed and the dependency graph was materialized to
\texttt{target/debug/deps}.

\paragraph{README.md}
Mathematical specification document covering: the Weaver Diff formal definition
$\Delta_W(S_1, S_2)$, the PI Conformance Residual $\mathcal{R}_{\Delta}$, the
Categorical Law $\Delta_W \neq \delta_P$, a worked JSON diff example showing the
two-tuple S1-to-S2 diff (Rename + Add), a Rust code listing of the full
\texttt{bridge.rs} implementation, and an explanation of how $\mathcal{R}_{\Delta} = 0$
is achieved by inserting \texttt{BridgeRx} before the conformance algorithm.

\subsection{Commands}
\label{sec:otel-weaver-exp-002-diff-to-residual-impl-commands}

The following commands are the operative build and test surface for EXP-002,
as inferred from the \texttt{Cargo.toml} and the evidence of materialized
\texttt{target/debug/deps} artifacts:

\begin{verbatim}
# Build the bridge library
cargo build

# Run the inline unit test
cargo test test_bridge_translation_equivalence

# Run the full wasm4pm integration suite
# (from /Users/sac/process-intelligence/sources/wasm4pm)
cargo test
\end{verbatim}

No \texttt{Makefile} or \texttt{cargo make} configuration is present in the
experiment directory. The experiment is a standalone library crate without a
binary target; it cannot be run directly but must be consumed by an upstream
integration harness. The \texttt{cargo test} output transcript for the standalone
experiment is not present in the repository; only the parent-scope
\texttt{GGEN\_OTEL\_WEAVER\_PI\_RUNTIME\_001.md} checkpoint contains a test
transcript (for the \texttt{wasm4pm} integration suite, not exp-002 in isolation).

\subsection{Data and Ontology Surfaces}
\label{sec:otel-weaver-exp-002-diff-to-residual-impl-data}

Two YAML mapping files at the \texttt{otel-weaver/mappings/} level constitute the
data surface governing EXP-002's boundary rules:

\paragraph{registry-diff-to-pi-drift-map.yaml}
Version 1.0.0 file declaring:
\begin{itemize}
    \item \texttt{nominal\_isolation\_declaration}: the statement that Weaver schema
          diffs must never be registered as process drift.
    \item Four \texttt{schema\_diff\_categories}: \texttt{schema\_addition},
          \texttt{attribute\_deprecation}, \texttt{type\_change}, and the implicit
          \texttt{attribute\_rename} case that \texttt{BridgeRx} handles. Each
          category declares \texttt{process\_drift\_impact: none} or
          \texttt{none\_unless\_unmapped}.
    \item Three \texttt{process\_drift\_definitions}: activity skip, sequence
          violation, and conformance score decay.
    \item Two \texttt{mapping\_boundary\_rules}: \texttt{drift\_isolation\_rule\_1}
          and \texttt{drift\_isolation\_rule\_2}.
\end{itemize}

\paragraph{otel-to-pi-witness-map.yaml}
Version 1.0.0 file aligned to OTel schema URL
\texttt{https://opentelemetry.io/schemas/1.25.0}. Declares three witness mapping
rules translating OTel service identity fields (\texttt{service.name},
\texttt{service.namespace}, \texttt{telemetry.sdk.name}) to PI feedstock keys
(\texttt{pi.feedstock.witness.id}, \texttt{pi.feedstock.source},
\texttt{pi.feedstock.payload\_size}). Includes a
\texttt{conformance\_isolation.prevention\_rule}: ``Never map feedstock directly
to court consequences.''

No TTL/RDF ontology surface exists for the S1/S2 schema types or the
\texttt{BridgeRx} translation semantics as of the current experiment status. This
is an open gap documented in
Section~\ref{sec:otel-weaver-exp-002-diff-to-residual-open}.

\subsection{Templates and Rendered Artifacts}
\label{sec:otel-weaver-exp-002-diff-to-residual-impl-templates}

EXP-002 is a pure Rust library crate; it does not use OTel Weaver's Jinja2/MiniJinja
template rendering capability. The template rendering pattern is used in exp-004
(registry-to-wasm4pm-compat-witness) to manufacture Rust witness type declarations
from the resolved registry JSON.

The experiment's relationship to Weaver templates is as their consumer: the output
of a Weaver template run (a set of Rust type declarations derived from the custom
PI registry defined in exp-001) would be the upstream schema surface that
\texttt{TelemetryFeedstockS1} and \texttt{TelemetryFeedstockS2} are designed to
parse. Whether exp-002's structs were manually authored or manufactured from a
Weaver template run is not documented in the repository.
\textsc{[AUTHOR THESIS / INTERPRETATION: the manual authorship of the S1/S2 structs
is an open question that EXP-004's generator pattern is designed to eventually
replace with registry-driven code manufacture.]}
